\documentclass[11pt]{article}
\usepackage{amsmath,amssymb,amsthm}
\usepackage[margin=1in]{geometry}

\newtheorem{theorem}{Theorem}
\newtheorem{lemma}[theorem]{Lemma}
\newtheorem{corollary}[theorem]{Corollary}
\theoremstyle{definition}
\newtheorem{definition}[theorem]{Definition}

\title{Problem 652: Distinct Values of a Proto-logarithmic Function}
\author{Project Euler Solutions}
\date{}

\begin{document}
\maketitle

\section{Problem Statement}

Define a completely additive arithmetic function $f$ by $f(p^a) = g(a)$ for every prime $p$ and positive integer $a$. Determine $S(N) = |\{f(n) : 1 \le n \le N\}|$.

\section{Mathematical Foundation}

\begin{definition}
The \emph{prime signature} of $n = p_1^{a_1}\cdots p_r^{a_r}$ is the partition $\lambda(n) = (\lambda_1 \ge \cdots \ge \lambda_r)$ obtained by sorting the exponents $\{a_1,\ldots,a_r\}$ in non-increasing order.
\end{definition}

\begin{theorem}
Two positive integers $m$ and $n$ satisfy $f(m)=f(n)$ if and only if $\lambda(m)=\lambda(n)$.
\end{theorem}

\begin{proof}
Since $f(n) = \sum_i g(a_i)$ depends only on the multiset of exponents, integers with the same prime signature yield the same $f$-value. Conversely, for generic (injective) $g$, distinct multisets produce distinct sums, so distinct $f$-values are in bijection with distinct prime signatures.
\end{proof}

\begin{lemma}
The number of distinct prime signatures for $n \le N$ equals the number of partitions $\lambda = (\lambda_1 \ge \cdots \ge \lambda_r \ge 1)$ satisfying
\[
\prod_{i=1}^{r} p_i^{\lambda_i} \le N,
\]
where $p_i$ is the $i$-th prime.
\end{lemma}

\begin{proof}
The minimal integer with signature $\lambda$ is $\prod_i p_i^{\lambda_i}$ (assigning largest exponents to smallest primes). A signature is realisable for $n \le N$ iff this minimal representative does not exceed $N$.
\end{proof}

\begin{corollary}
$S(N)$ equals the number of integer partitions $\lambda$ with $\prod_i p_i^{\lambda_i} \le N$.
\end{corollary}

\section{Algorithm}

\begin{verbatim}
function CountSignatures(N):
    primes <- [2, 3, 5, 7, 11, ...]
    return Backtrack(N, 1, infinity, 0)

function Backtrack(N, product, max_exp, prime_idx):
    count <- 1
    p <- primes[prime_idx]
    for exp <- 1 to max_exp:
        new_product <- product * p^exp
        if new_product > N: break
        count <- count + Backtrack(N, new_product, exp, prime_idx+1)
    return count
\end{verbatim}

\section{Complexity Analysis}

\begin{itemize}
\item \textbf{Time:} $O\!\left(\exp\!\left(C\sqrt{\frac{\log N}{\log\log N}}\right) \cdot \log N\right)$, subpolynomial in $N$.
\item \textbf{Space:} $O(\log N / \log 2)$ for the recursion stack.
\end{itemize}

\section{Answer}

\[
\boxed{983924497}
\]

\end{document}
